% Einheit 2 von 5 – Lineare Funktion f(x) = m·x + n (bank/lineare-funktionen/e2.jsonl, zusammenbau v0.1)
%% TODO Zweigzeile Teil 1: was hier gelernt wird (Satz in Schülersprache aus dem Einheitstitel) – Einheit 2
\einheitenkopf[e2]{Einheit 2 von 5 · Lineare Funktion f(x) = m·x + n}
\zweigzeile{ab Kl. 8 · P10 oft}
\setcounter{aufgabe}{11}

%% TODO Titel: Ich-kann-Satz fehlt in der Bank (Platzhalter: Kettenname) – Nr. 12
%% TODO Anweisung: ein Satz über der ersten Teilaufgabe fehlt in der Bank – Nr. 12
\begin{aufgabe}{m und n in der Gleichung markieren}
\begin{teile}
\teil m und n? Unterstreiche m, kreise n ein: $f(x) = 3x - 5$ m = \leerfeld , n = \leerfeld
\end{teile}
\end{aufgabe}

%% TODO Titel: Ich-kann-Satz fehlt in der Bank (Platzhalter: Kettenname) – Nr. 13
%% TODO Anweisung: ein Satz über der ersten Teilaufgabe fehlt in der Bank – Nr. 13
\begin{aufgabe}{Steigungsrichtung ohne Zeichnen}
\begin{teile}
\teil Steigt oder fällt die Gerade $f(x) = 3x - 4$? Kreuze an. \\ \kreuz{steigt} \\ \kreuz{fällt}
\end{teile}
\end{aufgabe}

%% TODO Titel: Ich-kann-Satz fehlt in der Bank (Platzhalter: Kettenname) – Nr. 14
%% TODO Anweisung: ein Satz über der ersten Teilaufgabe fehlt in der Bank – Nr. 14
\begin{aufgabe}{Steigungsdreieck lesen}
\swz{Einen Schritt nach rechts – wie viel nach oben oder unten? Lies am Steigungsdreieck ab.}{\begin{ksys}[xmin=-4,xmax=4,ymin=-5,ymax=5,ablesen] \gerade{3}{-1}{g} \steigungsdreieck{0}{-1}{3} \end{ksys}}
\end{aufgabe}

%% TODO \verfahren{…}: Verfahrensüberschrift in Sachsprache fehlt in der Bank (2.3 b)
%% TODO Titel: Ich-kann-Satz fehlt in der Bank (Platzhalter: Kettenname) – Nr. 15
%% TODO Anweisung: ein Satz über der ersten Teilaufgabe fehlt in der Bank – Nr. 15
\begin{aufgabe}{Graph zeichnen}
\swa{Wertetabelle: Fülle die Tabelle zu $f(x) = 3x - 2$ aus.}{\wertetabelle{x}{f(x)}{-1,0,1,2}}
\swa{Wertetabelle: Fülle die Tabelle zu $f(x) = 2x + 3$ aus.}{\wertetabelle{x}{f(x)}{-2,-1,0,1}}
\swa{Wertetabelle: Fülle die Tabelle zu $f(x) = -x + 5$ aus.}{\wertetabelle{x}{f(x)}{0,1,2,3}}
\swa{Wertetabelle: Fülle die Tabelle zu $f(x) = 4x - 3$ aus.}{\wertetabelle{x}{f(x)}{-1,0,1,2}}
\swa{Wertetabelle: Fülle die Tabelle zu $f(x) = -2x + 5$ aus.}{\wertetabelle{x}{f(x)}{0,1,2,3}}
\swfrage{Was bleibt gleich, was ändert sich?}
\swa{Punkte eintragen: Trage die Wertepaare $(-2|-2)$, $(-1|1)$, $(0|4)$, $(1|7)$ der Funktion $f(x) = 3x + 4$ ins Koordinatensystem ein.}{\begin{ksys}[xmin=-4,xmax=4,ymin=-4,ymax=8] \end{ksys}}
\swa{Gerade durch Punkte: Berechne zwei Punkte von $f(x) = 4x + 2$, trage sie ein und zeichne die Gerade durch sie.}{\begin{ksys}[xmin=-4,xmax=4,ymin=-4,ymax=7] \end{ksys}}
\swa{Mit n und Steigungsdreieck: Zeichne die Gerade zu $f(x) = 3x + 5$: Markiere n auf der y-Achse, gehe 1 nach rechts und m nach oben.}{\begin{ksys}[xmin=-4,xmax=4,ymin=-4,ymax=9] \end{ksys}}
\swa{Mit n und Steigungsdreieck: Zeichne die Gerade zu $f(x) = -3x + 4$.}{\begin{ksys}[xmin=-4,xmax=4,ymin=-4,ymax=5] \end{ksys}}
\swa{Mit n und Steigungsdreieck: Zeichne die Gerade zu $f(x) = \frac{1}{2}x + 3$.}{\begin{ksys}[xmin=-4,xmax=4,ymin=-4,ymax=5] \end{ksys}}
\swa{Mit n und Steigungsdreieck: Zeichne die Gerade zu $f(x) = 2x - 4$.}{\begin{ksys}[xmin=-4,xmax=4,ymin=-5,ymax=4] \end{ksys}}
\swa{Graph zeichnen: Zeichne den Graphen der Funktion f mit $f(x) = -4x + 3$ in das Koordinatensystem. \hfill (P10 2026 FOR)}{\begin{ksys}[xmin=-4,xmax=4,ymin=-4,ymax=4] \end{ksys}}
\end{aufgabe}

%% TODO \verfahren{…}: Verfahrensüberschrift in Sachsprache fehlt in der Bank (2.3 b)
%% TODO Titel: Ich-kann-Satz fehlt in der Bank (Platzhalter: Kettenname) – Nr. 16
%% TODO Anweisung: ein Satz über der ersten Teilaufgabe fehlt in der Bank – Nr. 16
\begin{aufgabe}{Ablesen}
\swz{n? Lies den y-Achsenabschnitt n der Geraden g ab.}{\begin{ksys}[xmin=-5,xmax=5,ymin=-5,ymax=5,ablesen] \gerade{1}{3}{g} \end{ksys}}
\swz{n? Lies den y-Achsenabschnitt n der Geraden g ab.}{\begin{ksys}[xmin=-5,xmax=5,ymin=-5,ymax=5,ablesen] \gerade{-2}{4}{g} \end{ksys}}
\swz{n? Lies den y-Achsenabschnitt n der Geraden g ab.}{\begin{ksys}[xmin=-5,xmax=5,ymin=-5,ymax=5,ablesen] \gerade{0.5}{-2}{g} \end{ksys}}
\swz{n? Lies den y-Achsenabschnitt n der Geraden g ab.}{\begin{ksys}[xmin=-5,xmax=5,ymin=-5,ymax=5,ablesen] \gerade{3}{-1}{g} \end{ksys}}
\swz{n? Lies den y-Achsenabschnitt n der Geraden g ab.}{\begin{ksys}[xmin=-5,xmax=5,ymin=-5,ymax=5,ablesen] \gerade{-1}{-4}{g} \end{ksys}}
\swfrage{Was bleibt gleich, was ändert sich?}
\swz{m? Lies die Steigung m der Geraden g mit einem Steigungsdreieck ab.}{\begin{ksys}[xmin=-5,xmax=5,ymin=-5,ymax=5,ablesen] \gerade{2}{-3}{g} \end{ksys}}
\swz{m? Lies die Steigung m der Geraden g mit einem Steigungsdreieck ab.}{\begin{ksys}[xmin=-5,xmax=5,ymin=-5,ymax=5,ablesen] \gerade{-2}{1}{g} \end{ksys}}
\swz{m? Lies die Steigung m der Geraden g ab. Gehe so weit nach rechts, bis du wieder auf einer Gitterlinie bist.}{\begin{ksys}[xmin=-5,xmax=5,ymin=-5,ymax=5,ablesen] \gerade{0.5}{1}{g} \end{ksys}}
\begin{teile}
\teil Welche Gleichung gehört zur Geraden? Kreuze an. \\ \kreuz{$f(x) = 3x - 4$} \\ \kreuz{$f(x) = -4x + 3$} \\ \kreuz{$f(x) = 3x + 4$} \\ \kreuz{$f(x) = -3x - 4$}
\end{teile}
\swa{Welche Gleichung gehört zu welcher Geraden? Ordne den Geraden f, g und h je eine der Gleichungen zu: $y = -3x + 3$; $y = -\frac{1}{3}x + 3$; $y = -3$; $y = 2x + 3$; $y = -3x$; $y = 2x - 3$. \hfill (P10 2019 OS) f: \leerfeld \quad g: \leerfeld \quad h: \leerfeld}{\begin{ksys}[xmin=-5,xmax=5,ymin=-5,ymax=5,ablesen] \gerade{-3}{3}{f} \gerade{2}{3}{g} \gerade{0}{-3}{h} \end{ksys}}
\end{aufgabe}

%% TODO Titel: Ich-kann-Satz fehlt in der Bank (Platzhalter: Kettenname) – Nr. 17
%% TODO Anweisung: ein Satz über der ersten Teilaufgabe fehlt in der Bank – Nr. 17
\begin{aufgabe}{Parameter deuten (steigend, fallend, parallel, Sonderfall m = 0)}
%% TODO Darstellung für schwach fehlt in der Bank (lineare-funktionen-e2-k6-s1-v1; Abschnitt „Für schwache Schüler“)
\swz[3]{Welche zwei Geraden sind parallel? $f(x) = 4x - 1$, $g(x) = -4x + 5$, $h(x) = 4x + 6$}{}
\end{aufgabe}

%% TODO Titel: Ich-kann-Satz fehlt in der Bank (Platzhalter: Kettenname) – Nr. 18
%% TODO Anweisung: ein Satz über der ersten Teilaufgabe fehlt in der Bank – Nr. 18
\begin{aufgabe}{Graph zeichnen – Fehler finden, Begründen, Darstellungswechsel}
\swz[3]{Ich finde den Fehler: Paul liest am Steigungsdreieck ab: 1 nach rechts, 3 nach oben, also $m = \frac{1}{3}$. Benenne den Fehler und gib m richtig an.}{\begin{ksys}[xmin=-4,xmax=4,ymin=-5,ymax=5,ablesen] \gerade{3}{-2}{g} \steigungsdreieck{0}{-2}{3} \end{ksys}}
\swfrage{Welche Gerade ist steiler: $f(x) = 3x - 7$ oder $g(x) = 5x + 4$? Begründe ohne Rechnung.}
\swz{Gleichung zur Geraden? Gib die Gleichung der Geraden g an.}{\begin{ksys}[xmin=-5,xmax=5,ymin=-5,ymax=5,ablesen] \gerade{-1}{3}{g} \end{ksys}}
\end{aufgabe}

\uebersichtskasten{Lineare Funktion: f(x) = m · x + n \\ m = Steigung: 1 nach rechts, m nach oben (m < 0: nach unten). \quad f(x) = 2x + 1: 1 nach rechts, 2 nach oben \\ n = y-Achsenabschnitt: die Gerade schneidet die y-Achse bei (0 | n). \quad f(x) = 2x + 1: bei (0 | 1)}
